\documentclass[10pt,a4paper]{amsart}
\usepackage[margin=19mm]{geometry}
\usepackage{amsmath,amssymb,mathtools}
\usepackage[scr=rsfs,cal=euler]{mathalfa}
\usepackage{xfrac}
\usepackage[unicode,hidelinks,bookmarks=false]{hyperref}
\newcommand{\ie}{i.e.\ }
\newcommand{\cf}{cf.\ }
\newcommand{\resp}{respectively\ }
\newcommand{\Subsection}{\subsection}
\providecommand{\nobreakdash}{\nobreak\hbox{-}}
\setlength{\emergencystretch}{2em}
\multlinegap0pt
\usepackage{amssymb}
\usepackage{tikz}
\newtheorem{prop}{Proposition}
\newtheorem{remark}{Remark}
\newtheorem{lemma}{Lemma}
\title{Function-Based Minimal Linear Codes over Galois Rings $\mathrm{GR}(p^{n}, \ell)$: Minimality Criteria and Infinite Constructions}
\author{Biplab Chatterjee, Sihem Mesnager, Ratnesh Kumar Mishra, Makhan Maji,, Kalyan Hansda}
\date{Source: arXiv:2603.26614v1 (CC BY 4.0)}
\begin{document}
\maketitle

\noindent\emph{Source and licence.} Function-Based Minimal Linear Codes over Galois Rings $\mathrm{GR}(p^{n}, \ell)$: Minimality Criteria and Infinite Constructions. Version arXiv:2603.26614v1 (CC BY 4.0), distributed by arXiv under the Creative Commons Attribution 4.0 International licence, http://creativecommons.org/licenses/by/4.0/, as recorded in the arXiv licence metadata for this exact version and shown on the abstract page https://arxiv.org/abs/2603.26614v1. Full licence notice: \texttt{licenses/arxiv-2603.26614-CC-BY-4.0.txt}.\par\smallskip

\noindent\textit{Editorial scope.} Counted unit: the article's lemma 26 (source0.tex lines 2517--2534). The article's complete original proof of it is delivered with it (source0.tex lines 2536--2604, one \texttt{proof} environment of 69 lines). No mathematics is written, repaired or re-derived: the statement and the proof are the article's own lines, in the article's own notation, and no retained line is substituted or re-flowed.  Retained uncounted as the source-local context the proof needs: lines 482--485; lines 498--568.  Nothing was omitted from any retained span, so the package carries no omission record.  No retained line contains a citation, so the wrapper carries no bibliography and no bibliography entry.  No retained line contains a figure environment, an image-inclusion command or drawing code, so no image is delivered and the package declares no asset dependency.  Editorial formatting only, in addition: an AMS \texttt{amsart} preamble that restates the article's own declarations and macros used by the retained lines, namely the environments \texttt{prop}, \texttt{remark}, \texttt{lemma} on separate counters, together with the standard packages those lines need.  The retained environments print in this document as Proposition 1, Remark 1, Lemma 1 in order of appearance, whereas the article prints the same items with the numbering of its own full document.  The complete native source of the article as distributed in the arXiv e-print for this exact version is bundled unchanged as \texttt{sources/197294/source0.tex}.
\begin{prop}\label{Proposition 11}
Let $v \in \mathrm{GR}(p^{n},\ell)^{m}$ be a nonzero vector that is not a root word.
Then $v^{\perp}$ is not free, and any minimal generating set of $v^{\perp}$ contains at least $m$ elements.
\end{prop}

\begin{remark}\label{Remark 12}
Let
\[
v' = p^{r}(u,v_2,\dots,v_m),
\qquad 1 \le r \le n-1,
\]
where $u$ is a unit in $\mathrm{GR}(p^{n},\ell)$.

Since $v'$ belongs to $p\,\mathrm{GR}(p^{n},\ell)^m$,
it is not a root word.
Consequently, by Proposition~\ref{Proposition 11},
its orthogonal module $v'^{\perp}$ is not free.

A minimal generating set of $v'^{\perp}$ is given by
\[
\left\{
(-v_2u^{-1},1,0,\dots,0),\dots,
(-v_mu^{-1},0,\dots,0,1),
(p^{\,n-r}u^{-1}-v_2u^{-1},1,0,\dots,0)
\right\}.
\]

Alternatively, one may replace the last generator and obtain
\[
\left\{
(-v_2u^{-1},1,0,\dots,0),\dots,
(-v_mu^{-1},0,\dots,0,1),
(p^{\,n-r}u^{-1},0,\dots,0)
\right\},
\]
which also forms a minimal generating set of $v'^{\perp}$.

More generally, for any unit $u'' \not\equiv 1 \pmod p$,
one may construct the set
\[
A=
\left\{
(p^{\,n-r}u^{-1}-v_2u^{-1},1,\dots,0),
\dots,
(p^{\,n-r}u^{-1}-v_mu^{-1},0,\dots,1),
(p^{\,n-r}u^{-1}-v_2u^{-1}u'',u'',0,\dots,0)
\right\},
\]
which again constitutes a minimal generating set of $v'^{\perp}$.

If we denote by
\[
A'=
\left\{
(p^{\,n-r}u^{-1}-v_2u^{-1},1,\dots,0),
\dots,
(p^{\,n-r}u^{-1}-v_mu^{-1},0,\dots,1)
\right\},
\]
then $A'$ is linearly independent, and its orthogonal complement
is the cyclic module generated by
\[
(u,-p^{\,n-r}+v_2,\dots,-p^{\,n-r}+v_m).
\]

Hence, $A$ provides a minimal generating set of
\[
\{p^{r}(u,v_2,\dots,v_m)\}^{\perp}
=
v'^{\perp}.
\]

This explicit description of orthogonal modules
for non-root vectors will be instrumental in the construction
of minimal codes in later sections.
\end{remark}

\begin{lemma}\label{Lemma 26}
Let
\[
v\in \mathrm{GR}(p^n,l)^m\setminus\{0\}
\]
be a root word. Then there exists a set
\[
\{\beta_1,\beta_2,\cdots,\beta_m\}
\subseteq O(p^{\,n-r}v),
\]
with $r\ge 1$, such that:
\[
1\leq w(\beta_i)\leq 2
\quad \text{and} \quad
v\cdot \beta_i=p^r
\]
for all $i$.
\end{lemma}

\begin{proof}
Since $v$ is a root word, at least one coordinate of $v$ is a unit.
Without loss of generality, write
\[
v=(u,v_2,v_3,\dots,v_m),
\]
where $u$ is a unit of $\mathrm{GR}(p^n,l)$.

By Proposition~\ref{Proposition 11}
and Remark~\ref{Remark 12},
a smallest generating set of $O(p^{\,n-r}v)$ is given by
\[
\Big\{
(-v_2u^{-1},1,0,\dots,0),\;
(-v_3u^{-1},0,1,\dots,0),\;
\dots,\;
(-v_mu^{-1},0,\dots,0,1),\;
(p^ru^{-1},0,\dots,0)
\Big\}.
\]

Each of these vectors has Hamming weight at most two.

From this generating set, consider instead the family
\[
\beta_i=
(p^ru^{-1}-v_iu^{-1},e_i)
\quad (2\le i\le m),
\]
together with
\[
\beta_1=(p^ru^{-1},0,\dots,0).
\]

Since linear combinations within a generating set preserve
the generated module, this new family also forms
a smallest generating set of $O(p^{\,n-r}v)$.

Now compute, for $2\le i\le m$,
\[
v\cdot \beta_i
=
u(p^ru^{-1}-v_iu^{-1})
+
v_i
=
p^r - v_i u^{-1}u + v_i
=
p^r.
\]

Similarly,
\[
v\cdot \beta_1
=
u(p^ru^{-1})
=
p^r.
\]

Thus,
\[
v\cdot \beta_i=p^r
\quad \text{for all } i,
\]
and clearly, each $\beta_i$ has weight 1 or 2.

This completes the proof.
\end{proof}

\end{document}
